\documentclass[11pt,a4paper]{article}
\usepackage[utf8]{inputenc}
\usepackage[T1]{fontenc}
\usepackage[margin=1in]{geometry}
\usepackage{amsmath}
\usepackage{amsfonts}
\usepackage{amssymb}
\usepackage{graphicx}
\usepackage{booktabs}
\usepackage{array}
\usepackage{xcolor}
\usepackage{hyperref}
\usepackage{listings}
\usepackage{tcolorbox}

\definecolor{atlasblue}{HTML}{3B3081}
\definecolor{atlasteal}{HTML}{42CCB2}
\definecolor{atlascoral}{HTML}{E26C48}

\hypersetup{
    colorlinks=true,
    linkcolor=atlasblue,
    urlcolor=atlasteal,
    pdftitle={Enterprise SDN Controller - Progress & Gap Analysis},
    pdfauthor={Mostafa Faouzi}
}

\title{\textbf{\Huge ENTERPRISE SDN CONTROLLER}\\
\vspace{0.3cm}
\large Implementation Progress \& Architecture Gap Analysis Report}
\author{\textbf{Mostafa Faouzi} \\
Master 2 RES, Sorbonne University \\
\vspace{0.2cm}
Prepared for: \textbf{Hamza ZERTA}}
\date{\today}

\begin{document}

\maketitle
\tableofcontents
\newpage

\section{Executive Summary}
This progress and gap analysis report details the implementation milestones achieved on the \textbf{Enterprise SDN Controller} orchestrator. The prototype has been successfully deployed on a dedicated GCP virtual machine hosting multi-vendor network simulation topologies (Nokia SR Linux and Dell OS10 virtual switches). 

All major REST endpoint services (38 endpoints), real-time gNMI topology stream bindings, dynamic configuration audits, and configuration drift-detection state transitions have been successfully engineered and verified. The code-level architecture is aligned with the Sorbonne University Master 2 RES consolidated technical blueprint. This document details the exact implementations completed, solved infrastructure bottlenecks, and provides a clear pathway for the final production deployment onto a SUSE RKE2 Kubernetes cluster.

\section{Implemented Architecture vs. Blueprint Reference}
The proposed production architecture relies on a converged 3-node SUSE RKE2 Kubernetes cluster orchestrating stateless microservice pods (Ingress, API gateway, Celery worker stream, telemetry collectors) paired with stateful cluster structures (HA PostgreSQL and Redis Sentinels). 

To ensure rapid prototyping and simplify debugging loops, the implementation was successfully executed on a single-node VM via a virtualized Docker-Compose environment simulating the components as dedicated containers.

\begin{table}[ht]
\centering
\caption{Architecture Mapping: Reference Blueprint vs. Current VM Deployment}
\vspace{0.2cm}
\begin{tabular}{p{3.5cm} p{5.5cm} p{5.5cm} c}
\toprule
\textbf{Layer} & \textbf{Production Blueprint} & \textbf{Current VM Implementation} & \textbf{Status} \\
\midrule
\textbf{Simulation} & PNETLab (20x Dell OS10 switches) & Containerlab (6x Nokia SRL, 2x Dell OS10) & \textcolor{green}{Active} \\
\textbf{Ingress} & MetalLB L2 VIP + NGINX Ingress & Nginx Docker container (TLS routing) & \textcolor{green}{Active} \\
\textbf{Stateless Apps} & FastAPI Gateway, gNMI Collector, Celery Workers (KEDA) & Monolithic FastAPI backend container, Fixed Celery worker & \textcolor{orange}{Merged} \\
\textbf{Stateful Tiers} & PostgreSQL HA (1P+2R), Redis Sentinel & PostgreSQL container (Single), Redis container (Single) & \textcolor{green}{Active} \\
\textbf{Storage Layer} & Longhorn CSI, MinIO S3 bucket & Docker local volumes \& bind mounts & \textcolor{green}{Active} \\
\textbf{Monitoring} & Rancher UI, Prometheus + Grafana & Flower Celery dashboard & \textcolor{orange}{Partial} \\
\bottomrule
\end{tabular}
\end{table}

\subsection{Architectural Alignment Rationale}
Operating a single-node Docker-Compose virtual network topology instead of a complete RKE2 cluster during active feature development has several advantages:
\begin{itemize}
    \item \textbf{Reduced Iteration Overhead:} Rebuilding and deploying updated code runs in under 15 seconds, compared to 3--5 minutes required to rebuild, push, pull, and cycle Kubernetes pods.
    \item \textbf{Simplified Diagnostic Surface:} Network tracing, database inspection, and process monitoring are co-located, isolating software bugs from cluster orchestration issues (e.g., PVC mount locks, service mesh routing mismatches).
\end{itemize}

\newpage
\section{Configuration Drift Detection \& Lifecycle Management}

The switch configuration lifecycle follows a state machine mapping:
\begin{center}
\texttt{DiscoveredRaw} $\rightarrow$ \texttt{CompliantActive} $\rightarrow$ \texttt{ConfigurationDrifted} $\rightarrow$ \texttt{TriggerRollback} $\rightarrow$ \texttt{CompliantActive}
\end{center}

\begin{figure}[ht]
\centering
\begin{tcolorbox}[colback=blue!5!white,colframe=atlasblue,title=State Transition Diagram]
\small
\begin{enumerate}
    \item \textbf{ZTP Onboarding:} Switch boots, queries DHCP option 67, runs the onboarding script, and registers with the controller (\texttt{DiscoveredRaw}).
    \item \textbf{Baseline Snapshot:} An approved initial configuration is generated and snapshot-saved in the database (\texttt{CompliantActive}).
    \item \textbf{Out-of-Band Mismatch:} An administrator modifies parameters (e.g., adding VLANs) directly on the switch CLI via console or SSH.
    \item \textbf{Drift State:} The background audit daemon detects a hash mismatch between the running config and baseline snapshot (\texttt{ConfigurationDrifted}).
    \item \textbf{Remediation:} The operator triggers a rollback to the baseline configuration, or clicks ``Accept Drift'' to promote the current live config to the new baseline.
\end{enumerate}
\end{tcolorbox}
\end{figure}

\subsection{Solved Infrastructure Challenges}
During system validation, two critical bugs were diagnosed and resolved to make drift detection functional:
\begin{itemize}
    \item \textbf{Transition from Console to SSH-based Collection:} \\
    Initially, the background discovery daemon queried Dell switches on the raw console socket (telnet on port \texttt{5000}). However, once the switch virtual machine completes its boot cycle inside the Containerlab container, Vrnetlab shuts down console communication. This caused connection timeouts in the collector. \\
    \textbf{Solution:} Rewrote the collector in \texttt{gnmi\_discovery.py} to connect via standard \textbf{SSH (Port 22)} using the robust \texttt{DellOS10Collector} driver class.
    
    \item \textbf{Command Echo \& Pager Pagination Mitigation:} \\
    When fetching the running configuration, the terminal pager triggered a \texttt{--More--} interactive prompt, causing incomplete configuration reads. Furthermore, command echo (e.g., \texttt{show running-configuration} at the first line of the output) caused text alignment mismatches against baseline snapshots. \\
    \textbf{Solution:} Configured the collector to send an escape command (\texttt{end}) followed by \texttt{terminal length 0} before config pulls. Added parsing logic in the collector to pop the echoed command, prune trailing terminal prompt lines, and strip carriage returns before saving configuration hashes to PostgreSQL.
\end{itemize}

\subsection{The ``Accept Drift'' Implementation}
We implemented an ``Accept Drift'' capability to complement the ``Rollback'' function. When a manual configuration change is validated, the operator can click ``Accept Drift as Baseline'' in the UI. 

This sends a POST request to \texttt{/api/v5/visibility/accept-drift}, creating a new \texttt{ConfigSnapshot} containing the live switch configuration, setting the new checksum hash as the switch baseline, and transitioning the switch state back to \texttt{compliant\_active}.

\newpage
\section{Verification and Test Coverage}
The platform includes an automated testing framework comprising:
\begin{itemize}
    \item \textbf{Unit \& Integration Tests:} 43 tests inside \texttt{test\_endpoints.py} verifying all 38 REST endpoints under a local SQLite database session manager.
    \item \textbf{Hybrid Discovery Tests:} SSH execution wrappers inside \texttt{test\_device\_inventory.py} validating live gNMI data structures.
\end{itemize}

\begin{lstlisting}[language=bash, caption=Verification Log of Switch Drift State inside PostgreSQL]
$ docker exec -t sdn_postgres_dev psql -U sdn_admin -d sdn_controller -c \
  "SELECT hostname, lifecycle_status, SUBSTRING(configuration_checksum, 1, 12) \
   FROM switches WHERE role = 'spine';"

 hostname | lifecycle_status | substring 
----------+------------------+-----------
 spine-01 | drifted          | 590df46167
 spine-02 | drifted          | a0e92f0166
(2 rows)
\end{lstlisting}

\section{Kubernetes Production Migration Path}
Upon feature completion, the transition from the single VM development environment to the target RKE2 production infrastructure is planned as follows:

\begin{enumerate}
    \item \textbf{Decouple Microservices:} Separate the FastAPI REST gateway, the Celery async worker, and the telemetry collectors into individual lightweight container images (defined in separate Dockerfiles).
    \item \textbf{Construct Kubernetes Deployment Manifests:} Author Helm charts wrapping resource allocations, Horizontal Pod Autoscalers (HPA), and KEDA scaled objects pointing to the Redis broker queue depth.
    \item \textbf{Provision Stateful Sets via Operators:}
    \begin{itemize}
        \item Deploy PostgreSQL High-Availability using the CloudNativePG Operator (1 primary, 2 replicas, configured with automated WAL archiving).
        \item Deploy Redis Sentinel clusters for task brokers.
    \end{itemize}
    \item \textbf{S3-Compatible Storage Integration (MinIO):} Update the python file storage logic inside \texttt{config\_lifecycle.py} to store configuration snapshot text strings inside S3 buckets hosted on MinIO using the \texttt{boto3} client library, offloading bulk text columns from PostgreSQL.
\end{enumerate}

\end{document}
